\section{Game Logic and Server Authority}
\label{sec:authority}

Turn control is held entirely on the board. \texttt{g->turn} is advanced only
after a move has been accepted and found not to end the game, so a rejected
request cannot hand the turn to the other player either.

A win is detected by testing the eight lines, three rows, three columns and two
diagonals, for three identical marks immediately after the cell is written; a
draw is the case where the board is full and no line was won. Both set
GAME\_OVER, after which only \texttt{NEW}, \texttt{STATE} and
\texttt{DISCONNECT} are accepted.

The authority claim is that \textbf{no rejected request can change state}, and
it rests on the structure quoted in Section~\ref{sec:impl} rather than on
inspection of every path: the three validation checks all return before the
sole write to the board. The claim is deliberately narrow enough to test from
the outside, and Section~\ref{sec:valid} describes how the suite asserts it by
reading the board with \texttt{STATE}, issuing a request expected to fail, and
reading it again.

The client is given no way to acquire authority. It holds no copy of the rules,
it derives whose turn it is only for display, and every redraw is driven by a
reply rather than by a local prediction. Replacing the client with a hostile one
therefore changes nothing about what positions are reachable: the worst a
modified client can do is send a request the server refuses.
